% 主表竖排，仅保留12个核心记号；局部记号在正文首次出现处定义。
\Needspace{95mm}
\section{符号说明}

\begin{table}[H]
\centering
\caption{主要符号与含义}
\label{tab:symbols}
\begin{tabularx}{0.82\linewidth}{@{}c>{\raggedright\arraybackslash}Xc@{}}
\toprule
符号 & 含义 & 单位 \\
\midrule
$\Omega$ & 目标区域 & -- \\
$p$ & 干扰源位置 & $\mathrm{m}$ \\
$x_t$ & 第 $t$ 个检测位置 & $\mathrm{m}$ \\
$f$ & 干扰源频道 & -- \\
$R_f$ & 频道 $f$ 的有效接收半径 & $\mathrm{m}$ \\
$\varepsilon$ & 示向度误差界 & $^\circ$ \\
$\widehat C_f$ & 频道 $f$ 的保守定位区域 & -- \\
$Q_f$ & 可靠接收测点区域 & -- \\
$z(C)$ & 区域 $C$ 的最小包围圆圆心 & $\mathrm{m}$ \\
$\rho(C)$ & 区域 $C$ 的最小包围圆半径 & $\mathrm{m}$ \\
$D$ & 定位区域直径 & $\mathrm{m}$ \\
$T$ & 任务耗时 & $\mathrm{s}$ \\
\bottomrule
\end{tabularx}
\end{table}
\par\medskip
